%======================================================================
\section{Introduction}\label{sec:intro}
%======================================================================

A group is sofic when its finite pieces can be modelled by permutations of finite sets, up to a small error. How large must those models be?
For a finite piece $E$ of a group, a \emph{chunk}, Cornulier's \emph{sofic profile} $\sigma_E(r)$ is the least number of points on which $E$ can be
modelled with multiplication correct on all but a fraction $1/r$ of the points and distinct elements separated almost everywhere~\cite{Cornulier2013}.
Amenable groups are sofic, so their profiles are finite. Cornulier asked whether a sofic group can have profile that is unbounded and not linear
in $r$~\cite[Problem~3.18]{Cornulier2013}. In~\cite{DouglasMghirbiA} we showed that Houghton's group $H_3$, which is finitely
presented~\cite{Brown1987}, with an explicit presentation due to Johnson~\cite{Johnson1999}, and elementary amenable, has a chunk whose profile is at least $\exp(\Omega(r^{c}))$, $c=1/(1+\log107)\approx0.129$. The proof rests on a counting criterion
(Theorem~\ref{thm:criterion}): if finitely many relations force many displaced copies of $S_3$ to commute cheaply, every model must be large,
because approximately commuting copies of $S_3$ carry entropy. That criterion can never give more than $\exp(O(r))$.

This paper shows that the criterion nearly reaches its ceiling in an amenable group, and that the group then becomes a quantum channel whose
repeated use needs nearly linear memory.

\begin{theoremA}
There is a finitely presented elementary amenable group $\Gamma$, with five generators, and a chunk $E$ of $\Gamma$ such that, for all $r\ge2$,
\[
\exp\Bigl(\frac{c\,r}{(\log r)^{12.75}}\Bigr)\ \le\ \sigma_E(r)\ \le\ \exp(C\,r\log r).
\]
So $\sigma_E(r)=\exp(r^{1+o(1)})$. The lower bound is certified by an explicit set of $25$ relations of $\Gamma$. Moreover, for every $\Delta\in(0,1]$ there are $c_\Delta,e_\Delta>0$ such that every assignment of unitaries in
$U(D)$ to the five generators under which the $25$ relations have defect at most $e\le e_\Delta$ in the normalized Hilbert--Schmidt norm, and
$ab$ stays at distance at least $\Delta$ from $I$, has $\log D\ge c_\Delta\,e^{-2}/(\log(1/e))^{25.49}$.
\end{theoremA}

Theorem~\ref{thm:main-group} gives the precise statement, and Proposition~\ref{prop:fp} the finite presentation, which follows from
Cornulier's criterion for permutational wreath products~\cite[Theorem~1.1]{Cornulier2006} once the affine group acting on the configurations is shown to be finitely presented. The lower bound is within a polylogarithmic factor in the exponent of the most the
criterion can give, and the upper bound shows that the truth is not much larger. The upper bound comes from models that hash configurations by
random affine maps; they are far smaller than the models that Følner sets give~\cite[Proposition~3.14]{Cornulier2013}, since the Følner function
of $\Gamma$ is doubly exponential (Proposition~\ref{prop:folner-lower}). The group $\Gamma$ also embeds in Brin's higher-dimensional Thompson group $3V$
(Proposition~\ref{prop:3V}), so $3V$, a finitely presented simple group whose soficity is open, has a chunk whose profile is finite and
$\exp(r^{1+o(1)})$.

\paragraph{Configuration lamps.} In Houghton's group the copies of $S_3$ sit on points of a ray; carrying a copy a distance $k$ takes a word of
length $k$, and commuting two copies at distance $k$ costs at least of order $k^2$ relations~\cite[Proposition~3.4]{DouglasMghirbiA}. A family of
$N$ copies spreads over distance $N$, and its commutators are expensive. We place the copies instead on the finitely
supported binary configurations of the three rays of $H_3$. Houghton's group moves the points, finitely supported affine maps act on the
configurations (a translation flips one point, a transvection adds one point into another), and one copy of $S_3$ sits on each configuration. The
$2^N$ configurations supported on the first $N$ points of a ray are reached by words of length $O(N^2)$, and two of them differ by a configuration
that at most $N$ transvections clear down to a single point. Each step touches at most three points, so its cost is governed by the Dehn function
of $H_3$ at scale $N$, which is polynomial by~\cite[Theorem~3.3]{DouglasMghirbiA}. The total cost of commuting two of the $2^N$ copies is therefore
polynomial in $N$: polylogarithmic in the number of copies. This is what the criterion needs.

\paragraph{A ladder.} The construction completes a sequence of transport mechanisms, each of which the criterion turns into a profile bound.

\begin{center}
\begin{tabular}{llll}
\toprule
transport & group & $\log\sigma_E(r)\ \gtrsim$ & profile finite \\
\midrule
along a ray & $H_3$~\cite{DouglasMghirbiA} & $r^{0.129}$ & yes \\
by one element & dyadic $\mathcal G$ (\S\ref{sec:dyadic}) & $r^{1/3}$ & yes \\
exponential orbits & $\mathcal L\le V$ (\S\ref{sec:thompson}) & $r/\log^2r$ & if $F$ is amenable \\
configurations & $\Gamma$ (\S\S\ref{sec:group}--\ref{sec:profile}) & $r/(\log r)^{12.75}$ & yes \\
\bottomrule
\end{tabular}
\end{center}

\noindent Transport by the powers of one element gives exponent $1/(p+1)$ when the far commutators have area of order $k^p$; along a ray of
$H_3$ the area is at least quadratic~\cite[Proposition~3.4]{DouglasMghirbiA}, so that transport stops at $1/3$, and $\mathcal G$ attains $1/3$ with
seventeen relations.
Thompson's group $F$ transports cheaply along exponentially growing orbits, which gives nearly exponential profile in a finitely presented
subgroup $\mathcal L$ of Thompson's group $V$, provided its models exist at all: they do if $F$ is amenable, and we show that $\mathcal L$ is
hyperlinear exactly when $F$ is, so that unitary models of this chunk exist whenever $F$ is hyperlinear
(Proposition~\ref{prop:thompson-hyperlinear}). The configuration lamps keep the exponential transport and drop the condition. For a variant on the
configurations of a grid (Remark~\ref{rem:grid}) a sketched argument gives a stronger lower bound, $r/(\log r)^{7.38}$, with the same upper bound.

\paragraph{Quantum memory.} Kretschmann and Werner showed that every causal quantum process is a concatenation of memory
channels~\cite{KretschmannWerner2005}, and Rybár and Ziman asked when a single memory device, which is not reset between uses, applies the same
channel to each of its inputs~\cite{RybarZiman2008}. In~\cite{DouglasMghirbiB} we asked how much memory such a device needs to apply the same quantum channel $n$ times, when
each output must be released before the next input arrives and every retained qubit counts. The device may export qubits and import fresh ones; the \emph{exchange}
$J$ counts the qubits exchanged, and the \emph{purity}, measured as the number of qubits minus the entropy of the initial memory and of the imports,
is a resource. The error is the half trace distance to the ideal experiment, against adaptive observers that keep quantum references. For a
gadget channel in the tracial closure, $\frac12\log r_*$ qubits per use, with $r_*$ the Choi rank, is the optimal exchange
rate~\cite[Propositions~3.1 and~9.1]{DouglasMghirbiB}. A channel built from a finite set of relations inherits the approximation profile of the
group: small memory would give small approximate representations. From $H_3$ we obtained an explicit channel needing memory $(n/\log n)^{0.069}$. The
configuration lamps give nearly linear memory, and a trade-off between memory and purity that is sharp up to polylogarithmic factors.

\begin{theoremB}
The certificate of Theorem~A defines an explicit quantum channel $\Phi$ on $\mathbb C^{873}$, in the closure of channels with finite tracial
factorizations. For causal services of $n$ uses of $\Phi$ with error at most $1/16$ and large $n$:
\begin{enumerate}[label=(\alph*),itemsep=1pt]
\item with purity at most $B$, where $\log n\le B\le c_0n$, the least memory lies between $c\,n/(B(\log n)^{25.49})$, at any exchange, and
$K(n/B)\log n$, attained by exact services that refresh their memory every round, with exchange $(n-1)Q$; so the least memory at purity $B$,
times $B$, is $n^{1+o(1)}$, and purity $O(\log n)$ forces memory $n^{1-o(1)}$;
\item at the optimal exchange rate, that is, with exchange $J\le\frac n2\log r_*$ over $n$ uses, the least memory lies between $c\sqrt n/(\log n)^{12.75}$, whatever the purity, and $K\sqrt{n\log n}$; services
near the least memory spend purity $n^{1/2\pm o(1)}$, at least that much by (a) and at most $3Q+6$, so at that rate memory and purity are both
about $\sqrt n$.
\end{enumerate}
\end{theoremB}

Theorem~\ref{thm:main-channel} gives the precise statement. Its engine is a one-round theorem for every gadget channel whose relations contain
those of $S_3$ and a name for $ab$ (Theorem~\ref{thm:weighted}): a single round of any device reveals an approximate representation of the group on the device's memory, with a
relator defect of the order of the square root of the leakage and entropy production of that round. In~\cite[Theorem~5.5 and
Appendix~B]{DouglasMghirbiB} this was proved for one gadget, built from Houghton's group. The upper bounds come from sofic models, which give devices directly (Proposition~\ref{prop:sofic-upper}): at most points of a model the
pure bath state reproduces the channel exactly, so a memory that is maximally mixed on those points serves the channel with no error at all. So
the memory of $\Phi$ is caught between two readings of the profile of $\Gamma$ at scale $n$, and for $\Gamma$ both are $n^{1+o(1)}$. At the
optimal exchange rate a device can spend at most $3Q+6$ bits of purity (Proposition~\ref{prop:purity-budget}), which gives the lower bound in
part~(b). The upper bound needs a device that keeps a sofic model in memory for all $n$ rounds instead of refreshing it: a streaming theorem
(Theorem~\ref{thm:streaming}) uses the good points of a model as an exact dilation whose bath grows slightly in every round, so that the bad points
cost purity rather than accumulated error.

\paragraph{Physics.} A local lattice system whose bath starts in its maximally mixed state, run for time $T$, produces channels that are cheap
for fixed $T$: they need memory only polynomial in
$T$ and $\log n$ (Proposition~\ref{prop:light-cone}). So a local process that produces $\Phi$ to accuracy $\delta$ takes time at least about
$\delta^{-1/\nu}$ in lattice dimension $\nu$ (Corollary~\ref{cor:preparation}). Autonomous dynamics with a bath in its tracial state face a different constraint: by the mean ergodic
theorem, their long-time limits are compressions of conditional expectations and hence self-adjoint
channels~\cite[Remark~10.3]{JungeLeMerdyXu2006},~\cite[Remark~5.2]{HaagerupMusat2011}, and they satisfy a further inequality between coherence and return (Propositions~\ref{prop:autonomous}
and~\ref{prop:return}). Every channel $\Psi\Psi^*$ with $\Psi$ in the tracial closure is such a limit (Proposition~\ref{prop:nonlocal}), but for gadget channels this route gives nothing
expensive: composing a gadget channel with its adjoint erases the multiplication of the group, and $\Phi\Phi^*$ has an exact finite factorization
on at most $3r_*-2$ qubits (Proposition~\ref{prop:compositions}). Averaging the doubling of $\Phi$ with the identity does give an expensive
limit: a bounded time-independent Hamiltonian coupled to a tracial bath produces a channel on $M_{1746}$ exactly at the times $\frac\pi2k$, and in
the limit of long times, that trades memory and purity at the same rate as $\Phi$ (Proposition~\ref{prop:lazy}). The Hamiltonian is not local;
whether local autonomous dynamics can produce an expensive channel is open.

\paragraph{Dependencies on the companion papers.} Table~\ref{tab:dependencies} lists every result of~\cite{DouglasMghirbiA}
and~\cite{DouglasMghirbiB} used here, with the numbering of their version 1.1.

\begin{table}[ht]
\centering\small
\begin{tabular}{@{}p{0.47\textwidth}p{0.2\textwidth}p{0.25\textwidth}@{}}
\toprule
result & source & used in\\
\midrule
counting criterion & \cite[Thm.~8.1]{DouglasMghirbiA} & Thm.~\ref{thm:criterion}, every lower bound\\
transport by the powers of one element & \cite[Cor.~8.2]{DouglasMghirbiA} & Section~\ref{sec:criterion}\\
Johnson's presentation of $H_3$ and its conventions & \cite[Section~2]{DouglasMghirbiA} & Section~\ref{sec:houghton}\\
far-commutator area and polynomial Dehn bound for $H_3$ & \cite[Thms.~3.2, 3.3]{DouglasMghirbiA} & Prop.~\ref{prop:dehn}\\
quadratic lower bound along a ray & \cite[Prop.~3.4]{DouglasMghirbiA} & Sections~\ref{sec:intro}, \ref{sec:criterion}, \ref{sec:ladder}\\
bounded or at least $2r$ (sharpening Cornulier's Fact~3.5); embeddings in finite groups & \cite[Lemma~2.1]{DouglasMghirbiA} & Section~\ref{sec:ladder}, after Thm.~\ref{thm:main-group}\\
superpolynomial profile of $H_3$ & \cite[Thm.~1.1]{DouglasMghirbiA} & Section~\ref{sec:thompson}\\
projection products & \cite[Lemma~4.1]{DouglasMghirbiA} & Lemma~\ref{lem:projection-product}\\
\midrule
causal services, exchange, purity and error & \cite[Def.~2.3, Section~3]{DouglasMghirbiB} & Section~\ref{sec:channel}\\
gadget channels: Choi rank $r_*$, flat probe, entropy $\log r_*$ & \cite[Prop.~9.1]{DouglasMghirbiB} & Section~\ref{sec:channel-setting}\\
the rank rate is the optimal exchange rate of gadget channels in $\cK_d$ & \cite[Props.~3.1, 9.1]{DouglasMghirbiB} & Thm.~B(b)\\
sequential extraction & \cite[Thm.~5.2]{DouglasMghirbiB} & Thm.~\ref{thm:extraction}\\
purity budget $14Q+3$ at the rank rate (sharpened here) & \cite[Thm.~5.8]{DouglasMghirbiB} & Prop.~\ref{prop:purity-budget}\\
entropy variance & \cite[Lemma~5.9]{DouglasMghirbiB} & Lemma~\ref{lem:surprisal}\\
rolling-reservoir compiler & \cite[Appendix~A, Lemmas~A.1, A.4]{DouglasMghirbiB} & Thm.~\ref{thm:streaming}\\
channel profiles are group profiles; hyperlinear groups give channels in $\cK_d$ & \cite[Thm.~9.2]{DouglasMghirbiB} & Section~\ref{sec:channel}\\
one-round theorem for the Houghton gadget (generalised here) & \cite[Thm.~5.5, Appendix~B]{DouglasMghirbiB} & Thm.~\ref{thm:weighted}, Section~\ref{sec:weighted-proof}\\
weighted norms; charged transport and weighted extraction for the Houghton gadget (generalised here) & \cite[Appendix~B, Lemmas~B.2, B.4]{DouglasMghirbiB} & Lemmas~\ref{lem:weighted-extraction}, \ref{lem:charged-area}\\
weighted $S_3$ rounding & \cite[Lemma~B.3]{DouglasMghirbiB} & Lemma~\ref{lem:rounding}\\
weighted sector bound (proved here by sequential measurement) & \cite[Thm.~B.1]{DouglasMghirbiB} & Thm.~\ref{thm:sector}\\
\bottomrule
\end{tabular}
\caption{Results of the companion papers used in this paper (numbering of their version 1.1).}
\label{tab:dependencies}
\end{table}

\paragraph{Organization.} Section~\ref{sec:criterion} recalls the counting criterion. Section~\ref{sec:ladder} treats the dyadic group and the
subgroup of Thompson's group $V$. Sections~\ref{sec:group}--\ref{sec:profile} construct the configuration lamps, show that they are finitely presented and
embed in $3V$, prove the local fillings and the pair-area theorem, and prove Theorem~A. Section~\ref{sec:channel} proves Theorem~B, with the one-round theorem proved in
Section~\ref{sec:weighted-proof}. Section~\ref{sec:physics} treats local physics, and Section~\ref{sec:open} lists open problems. The certificate of
$25$ relations is in the accompanying files, together with programs that check it; logarithms are to base two unless written $\ln$.
